\section{Methodology}

\subsection{Framework Architecture}

% ~120 words — 6-stage pipeline, 24kHz, model-agnostic
The framework implements a six-stage pipeline
(Figure~\ref{fig:architecture}): (1)~dataset preparation, including
standard corpora and our novel challenge set; (2)~speech generation through
a model-agnostic client interface, where each TTS system inherits from a
common \texttt{BaseTTSClient} class configured via YAML, with exponential
backoff for API rate limits; (3)~per-utterance metric computation across 28
objective measures; (4)~normalization and composite score aggregation;
(5)~visualization; and (6)~statistical significance testing. All generated
audio is resampled to 24\,kHz mono float32 to ensure fair comparison across
systems with differing native sample rates. The modular design allows new
models and metrics to be added by implementing the client interface and
registering metric functions, without modifying the core pipeline.

\subsection{Evaluation Metrics}

% ~200 words — 28 metrics, 6 dimensions, table + prose descriptions
We evaluate synthesized speech across 28 metrics organized into six
perceptual dimensions (Table~\ref{tab:dimensions}).

\begin{table}[t]
\centering
\caption{Evaluation dimensions and their constituent metrics. All 28 metrics
are computed per utterance and aggregated per dimension.}
\label{tab:dimensions}
\small
\begin{tabular}{@{}p{1.6cm}p{4.4cm}c@{}}
\toprule
\textbf{Dimension} & \textbf{Metrics} & \textbf{N} \\
\midrule
Naturalness & UTMOS, DNSMOS, output SNR, PESQ, STOI & 5 \\
Intelligibility & WER, CER, ASR mismatch, word skip rate, semantic distance & 5 \\
Speaker Sim. & Resemblyzer cosine similarity & 1 \\
Prosody & F0 range, jitter, shimmer, HNR, pause ratio, speaking/syllable rate, pause count/mean dur., duration ratio, energy mean/std, dynamic range & 13 \\
Robustness & repetition, silence anomaly, insertion rate & 3 \\
Latency & real-time factor (RTF) & 1 \\
\bottomrule
\end{tabular}
\end{table}

\textbf{Naturalness} combines neural MOS prediction via
UTMOS~\cite{saeki2022utmos} and DNSMOS~\cite{reddy2021dnsmos} with
signal-level fidelity measured by PESQ~\cite{itu2001pesq} and
STOI~\cite{taal2011stoi}, plus VAD-based output SNR as a reference-free
signal quality indicator.
\textbf{Intelligibility} computes word and character error rates using
Whisper base~\cite{radford2023whisper} as the ASR backbone, supplemented by
word skip rate (proportion of deleted words) and round-trip semantic
distance via sentence embeddings to capture meaning preservation beyond
surface-level transcription accuracy.
\textbf{Speaker similarity} measures cosine distance between reference and
synthesized speaker embeddings using
Resemblyzer~\cite{jemine2019resemblyzer}; it is evaluated only when
reference audio is available.
\textbf{Prosody} captures 13 features spanning pitch dynamics (F0 range,
jitter, shimmer, HNR), temporal patterns (pause ratio, speaking and
syllable rates, pause count and mean duration, duration ratio), and energy
statistics (mean, standard deviation, dynamic range)---all extracted via
Parselmouth~\cite{jadoul2018parselmouth} and Silero VAD.
\textbf{Robustness} flags three failure modes: n-gram repetition (the same
phrase repeated three or more times), silence anomalies from VAD analysis,
and spurious word insertions.
\textbf{Latency} records the real-time factor (RTF) from generation
metadata, measuring computational cost relative to audio duration.

\subsection{Challenge Set}

% ~170 words — 7 categories, ~350 utterances, code-switching theme
Standard TTS benchmarks rely on clean read speech such as
Seed-TTS-Eval~\cite{anastassiou2024seedtts}, which fails to expose
system-specific failure modes. We introduce a 7-category adversarial
challenge set comprising approximately 350 utterances (${\sim}50$ per
category), each targeting a distinct synthesis difficulty:
\textit{numbers} (dates, currency, phone numbers), \textit{proper nouns}
(foreign-origin names, brands), \textit{tongue twisters} (phonetically
demanding sequences), \textit{long-form} passages (200+ words testing
coherence over extended output), \textit{questions} (interrogative prosody
requiring rising intonation), \textit{emotional} text (excitement, sadness,
anger), and \textit{code-switching} (English with embedded French and
German phrases). The code-switching category directly addresses multilingual
speech production, aligning with the Interspeech 2026 ``Speaking Together''
theme. Concurrent work on adversarial TTS
evaluation~\cite{manku2025emergentttseval} introduces 1,645 challenge cases
but relies on LLM-as-judge scoring; our challenge set instead pairs each
utterance with ground-truth transcripts and evaluates all 28 objective
metrics, enabling reproducible, fine-grained difficulty analysis across
categories and systems.

\subsection{Use-Case Composite Scoring}

% ~180 words — normalization, weights table, renormalization, Wilcoxon
Raw dimension scores are derived in three steps. First, each metric is
clipped to empirically determined bounds and linearly scaled to $[0,\,1]$.
Second, lower-is-better metrics (e.g., WER, RTF) are inverted so that
higher always indicates better performance. Third, prosodic measures where
both extremes are undesirable---jitter, shimmer, speaking rate, pause
ratio, syllable rate, and duration ratio---use target-based deviation
scoring: values near the empirical optimum receive scores close to 1.0,
with symmetric or asymmetric penalties for deviation. Each dimension score
is the mean of its constituent normalized metrics.

Use-case composites are weighted sums of the six dimension scores
(Table~\ref{tab:composites}), reflecting deployment-specific priorities:
audiobook narration emphasizes naturalness and prosody, while conversational
AI prioritizes latency and intelligibility. When a dimension cannot be
computed (e.g., speaker similarity without reference audio), the remaining
weights are renormalized by dividing by their sum, preserving a $[0,\,1]$
range.

\begin{table}[t]
\centering
\caption{Use-case composite weights. Each row sums to 1.0; ``---'' denotes
zero weight (dimension excluded from scoring).}
\label{tab:composites}
\small
\begin{tabular}{@{}lcccccc@{}}
\toprule
\textbf{Use Case} & \textbf{Nat} & \textbf{Int} & \textbf{Spk} & \textbf{Pro} & \textbf{Rob} & \textbf{Lat} \\
\midrule
Audiobook      & .35 & .15 & .10 & .25 & .15 & --- \\
Conversational & .20 & .25 & --- & .10 & .15 & .30 \\
Voice Cloning  & .20 & .15 & .40 & .05 & .10 & .10 \\
Low-Latency    & .15 & .20 & --- & .05 & .15 & .45 \\
Balanced       & .20 & .20 & .15 & .10 & .15 & .20 \\
\bottomrule
\end{tabular}
\end{table}

Statistical significance of pairwise model differences is assessed via the
Wilcoxon signed-rank test on matched per-utterance scores, with Bonferroni
correction for multiple comparisons ($\alpha = 0.05 / \binom{k}{2}$ for $k$
models). Tests are applied to both the balanced composite and raw UTMOS
scores, providing complementary views of overall and perceptual quality
differences.
